%% 5-7 いちばん高いのはどこ？：最大と最小（第11回：書き足し。ページ数は自由・偶数）
\bkchapter{5}{7}{いちばん高いのはどこ？\\最大と最小}{\bklevel{3}}{いちばん高いのはどこ？　最大と最小}
\bkmichishirube{2-2,3-10,5-5,5-6}{5-8,6-13}

\begin{goal}
\begin{itemize}[label={\color{bkTheme}●}, leftmargin=1.6\zw, itemsep=0pt, topsep=1pt]
\item \bklevel{3}\enspace 導関数の符号から、グラフが上がるか下がるかを判断できる。
\item \bklevel{3}\enspace 増減表で、山（極大）と谷（極小）を見つけられる。
\item \bklevel{3}\enspace 箱の体積や面積がいちばん大きくなる形を、微分で求められる。
\end{itemize}
\end{goal}

\begin{hatena}
\noindent\begin{minipage}[c]{0.6\linewidth}\raggedright
1辺\qty{12}{cm}の正方形の紙の4すみから、同じ大きさの正方形を切りとって折り曲げ、ふたのない箱をつくります。

切りとる正方形を小さくすると、箱は浅くなります。大きくすると、底がせまくなります。体積がいちばん大きくなるのは、何 cm の正方形を切りとったときでしょう。
\end{minipage}\hfill
\begin{minipage}[c]{0.37\linewidth}\centering
\begin{tikzpicture}[x=1mm, y=1mm, scale=0.62]
  \fill[bkYellow!20] (0,0) rectangle (36,36);
  \foreach \cx/\cy in {0/0,27/0,0/27,27/27} \fill[bkGray!40] (\cx,\cy) rectangle ++(9,9);
  \draw[bkInk, line width=0.8pt] (0,0) rectangle (36,36);
  \draw[bkInk, line width=0.5pt, dashed] (9,0) -- (9,36) (27,0) -- (27,36) (0,9) -- (36,9) (0,27) -- (36,27);
  \node[font=\sffamily\scriptsize, anchor=north] at (18,-0.5) {\qty{12}{cm}};
  \node[font=\sffamily\scriptsize] at (4.5,4.5) {$x$};
  \draw[bk arrow, bkGray, line width=0.9pt] (39,18) -- (45,18);
  \begin{scope}[shift={(48,6)}, x={(1mm,0mm)}, y={(0.5mm,0.35mm)}, z={(0mm,1mm)}]
    \draw[bkInk, line width=0.7pt, fill=bkYellow!35] (0,0,0) -- (18,0,0) -- (18,0,9) -- (0,0,9) -- cycle;
    \draw[bkInk, line width=0.7pt, fill=bkYellow!20] (18,0,0) -- (18,18,0) -- (18,18,9) -- (18,0,9) -- cycle;
    \draw[bkInk, line width=0.7pt] (0,0,9) -- (0,18,9) -- (18,18,9);
  \end{scope}
\end{tikzpicture}
\end{minipage}
\end{hatena}

いちばん大きい所や、いちばん小さい所を見つけることは、物をつくるときにいつも必要になります。材料を少なくしたい、たくさん入れたい、いちばん速く着きたい……。この章では、微分を使って、こうした「いちばん」を見つけます。

\clearpage
%% ---------------- やってみよう・思い出そう ----------------
\begin{yattemiyou}[箱をつくって体積をくらべよう]
\par\addvspace{3pt}\textbf{用意するもの}\par\noindent 方眼紙（1辺\qty{12}{cm}の正方形を5まい）、はさみ、テープ
\begin{enumerate}[label={\bkcircled[bkGreen]{\arabic*}}, leftmargin=2.4\zw, itemsep=1pt, topsep=3pt]
\item 切りとる正方形の1辺 $x$ を1、2、3、4、\qty{5}{cm}にして、5つの箱をつくろう。
\item 底の1辺・深さをはかり、体積（底の面積×深さ）を計算して、表にまとめよう。
\item どの箱の体積がいちばん大きいかな？
\end{enumerate}
\end{yattemiyou}

\noindent\begin{minipage}[c]{0.5\linewidth}\raggedright
\textbf{底の1辺は $12-2x$}\quad 右の図のように、紙の1辺\qty{12}{cm}から、左はしで $x$、右はしでも $x$ を切りとります。残りが箱の底の1辺なので、$12-x-x=12-2x$ です。深さは、折り曲げた部分の $x$ です。
\end{minipage}\hfill
\begin{minipage}[c]{0.46\linewidth}\centering
\begin{tikzpicture}[x=1mm, y=1mm]
  \fill[bkYellow!20] (0,0) rectangle (48,48);
  \fill[bkBlue!18] (12,12) rectangle (36,36);
  \foreach \cx/\cy in {0/0,36/0,0/36,36/36} \fill[bkGray!40] (\cx,\cy) rectangle ++(12,12);
  \draw[bkInk, line width=0.8pt] (0,0) rectangle (48,48);
  \draw[bkInk, line width=0.5pt, dashed] (12,0) -- (12,48) (36,0) -- (36,48) (0,12) -- (48,12) (0,36) -- (48,36);
  \draw[bkInk, line width=0.5pt, |<->|] (0,-3) -- (12,-3);
  \draw[bkInk, line width=0.5pt, |<->|] (12,-3) -- (36,-3);
  \draw[bkInk, line width=0.5pt, |<->|] (36,-3) -- (48,-3);
  \node[font=\sffamily\footnotesize, anchor=north] at (6,-4.5) {$x$};
  \node[font=\sffamily\footnotesize, anchor=north] at (24,-4.5) {$12-2x$};
  \node[font=\sffamily\footnotesize, anchor=north] at (42,-4.5) {$x$};
  \node[font=\sffamily\footnotesize, align=center] at (24,24) {箱の底};
  \node[font=\sffamily\scriptsize, align=center] at (24,42) {折って\\かべに};
\end{tikzpicture}
\end{minipage}

\medskip
体積は「底の面積×深さ」なので、
\[ V=(12-2x)\times(12-2x)\times x=x(12-2x)^2 \]
です。$x=2$ なら、底の1辺は $12-4=8$、体積は $8\times8\times2=128$ です。

\begin{bkcenter}
\begin{bktable}{colspec={Q[l,m,wd=34mm] *{5}{Q[c,m,wd=12mm]}}, rowsep=1.5pt}
$x$（cm） & 1 & 2 & 3 & 4 & 5 \\
底の1辺 $12-2x$（cm） & 10 & 8 & 6 & 4 & 2 \\
体積 $V$（\unit{cm^3}） & 100 & 128 & 108 & 64 & 20 \\
\end{bktable}
\bktabcaption{切りとる正方形の1辺と、箱の体積}
\end{bkcenter}

5つの中では、$x=2$ のときがいちばん大きくなりました。でも、$x=1.9$ や $x=2.1$ のほうが大きいかもしれません。電卓で計算すると、$x=1.9$ で $127.756$、$x=2.1$ で $127.764$ で、どちらも128より少しだけ小さくなりました。では、$x=2.01$ では？　表をいくら細かくしても、「ほんとうにいちばん」とは言いきれません。微分を使うと、これをはっきり決められます。

\noindent\textbf{$x$ が使える範囲}\quad 切りとる正方形の1辺 $x$ は0より大きい数です。また、底の1辺 $12-2x$ も0より大きくないと、箱になりません。$12-2x>0$ より $x<6$。だから、$x$ は $0<x<6$ の範囲の数です。

\begin{omoidasou}
\bkrecall{5-3}{$f'(a)$ は、$x=a$ での接線の傾き。$f'(a)=\displaystyle\lim_{h\to0}\frac{f(a+h)-f(a)}{h}$。}
\bkrecall{5-4}{$(x^n)'=nx^{n-1}$。定数倍・和は1つずつ微分する。}
\bkrecall{5-5}{$f'(x)$ をもう一度微分したものが $f''(x)$（2回微分）。位置を2回微分すると加速度。}
\bkrecall{5-6}{積の微分 $(fg)'=f'g+fg'$。$\{(ax+b)^n\}'=n(ax+b)^{n-1}\times a$。}
\bkrecall{2-2}{$a>0$ は正、$a<0$ は負。$(+)\times(-)=(-)$、$(-)\times(-)=(+)$。}
\bkrecall{3-8}{$(a-b)^2=a^2-2ab+b^2$。たとえば $(12-2x)^2=144-48x+4x^2$。}
\bkrecall{3-10}{和が $-8$、積が12になる2つの数は $-2$ と $-6$ なので、\\$x^2-8x+12=(x-2)(x-6)$。$=0$ になる $x$ は2と6。また、$x^2-1=(x+1)(x-1)$。}
\end{omoidasou}

%% ---------------- 数学の道具 ----------------
\Needspace{8\baselineskip}
\bksection{doguu}{数学の道具}

\bksubsection{1}{導関数の符号で、上がる・下がるがわかる}

\label{bk:5-7:fugou}$f(x)=x^3-3x$ のグラフで考えましょう。グラフの上のいろいろな点に、接線を引いてみます。

\noindent\begin{minipage}[c]{0.5\linewidth}\centering
\begin{tikzpicture}
\begin{axis}[bkplot, width=62mm, height=54mm, xmin=-2.3, xmax=2.5, ymin=-3, ymax=3.6, xtick={-1,1}, ytick={-2,2}, xlabel={$x$}, ylabel={$y$}]
  \addplot[bkcurve, domain=-2.1:2.1] {x^3-3*x};
  \foreach \a in {-1.7,1.7} {\edef\tmp{\noexpand\addplot[bkGreen!80!black, line width=1.1pt, domain=\a-0.18:\a+0.18] {\a^3-3*\a+(3*\a*\a-3)*(x-\a)};}\tmp}
  \foreach \a in {-0.5,0.5} {\edef\tmp{\noexpand\addplot[bkRed, line width=1.1pt, domain=\a-0.4:\a+0.4] {\a^3-3*\a+(3*\a*\a-3)*(x-\a)};}\tmp}
  \node[font=\sffamily\scriptsize, text=bkGreen!60!black, anchor=east] at (axis cs:-1.95,-1.3) {傾き＋};
  \node[font=\sffamily\scriptsize, text=bkRed!80!black, anchor=south west] at (axis cs:0.3,0.6) {傾き－};
  \node[font=\sffamily\scriptsize, text=bkGreen!60!black, anchor=west] at (axis cs:1.95,-1.3) {傾き＋};
\end{axis}
\end{tikzpicture}
\end{minipage}\hfill
\begin{minipage}[c]{0.46\linewidth}\raggedright
グラフが右上がりの所（緑）では、接線も右上がりで、傾きは正です。右下がりの所（赤）では、接線も右下がりで、傾きは負です。接線の傾きは $f'(x)$ でした（5-3）。
\end{minipage}

\medskip
だから、次のことが言えます。
\begin{itemize}[leftmargin=1.6\zw, itemsep=1pt, topsep=2pt, label={\color{bkTheme}●}]
\item $f'(x)>0$ の所では、$x$ がふえると $f(x)$ もふえる（\textbf{増加}）。
\item $f'(x)<0$ の所では、$x$ がふえると $f(x)$ はへる（\textbf{減少}）。
\end{itemize}
\bkidx{ぞうか・げんしょう}{増加・減少}

\noindent\textbf{数で確かめる}\quad $f'(x)=3x^2-3$ です。
\begin{itemize}[leftmargin=1.6\zw, itemsep=1pt, topsep=2pt, label={\color{bkTheme}●}]
\item $x=1.7$ では $f'(1.7)=3\times2.89-3=5.67>0$。\\$x$ を少しふやすと、$f(1.7)=-0.187$ → $f(1.8)=0.432$ で、ふえている。
\item $x=0.5$ では $f'(0.5)=0.75-3=-2.25<0$。\\$x$ を少しふやすと、$f(0.5)=-1.375$ → $f(0.6)=-1.584$ で、へっている。
\end{itemize}

\noindent\textbf{なぜそう言えるのか}\quad $f'(a)$ は、$\dfrac{f(a+h)-f(a)}{h}$ の、$h$ を0に近づけたときの近づく先でした（5-3）。$f'(a)$ が正なら、$h$ が小さい正の数のとき、この分数も正です。分母の $h$ が正なので、分子 $f(a+h)-f(a)$ も正、つまり $f(a+h)>f(a)$ です。少し右へ行くと、値が大きくなるのです。

\noindent\textbf{例}\quad 投げ上げたボールの高さ $x=20t-5t^2$ なら、$v=\dot{x}=20-10t$（5-5）。$t<2$ では $v>0$ で上がり、$t>2$ では $v<0$ で下がります。上がるのから下がるのに変わる $t=2$（$v=0$）で、いちばん高くなります。

\bksubsection{2}{増減表}

\label{bk:5-7:zougen}$f(x)=x^3-3x$ が、どこで上がり、どこで下がるかを、表にまとめましょう。次の順に調べます。\bkidx{ぞうげんひょう}{増減表}

\noindent\textbf{手順1　微分して、$f'(x)=0$ となる $x$ を求める}
\[ f'(x)=3x^2-3=3(x^2-1)=3(x+1)(x-1) \]
（3でくくり、$x^2-1=(x+1)(x-1)$ と因数分解した。3-10）。$f'(x)=0$ になるのは、$x+1=0$ か $x-1=0$ のときなので、$x=-1$ と $x=1$ です。

\noindent\textbf{手順2　$f'(x)$ の符号を調べる}\quad $f'(x)$ は、3と $(x+1)$ と $(x-1)$ のかけ算です。それぞれの符号を表にすると、かけ算の符号がわかります（2-2）。

\begin{bkcenter}
\begin{tblr}{colspec={Q[c,m,wd=15mm] Q[c,m,wd=18mm] Q[c,m,wd=12mm] Q[c,m,wd=24mm] Q[c,m,wd=12mm] Q[c,m,wd=18mm]}, colsep=4pt, hlines={0.6pt,bkLine}, vline{2}={0.6pt,bkLine}, rowsep=2pt, row{1}={bg=bkTheme!10, font=\small}}
$x$ & $x<-1$ & $-1$ & $-1<x<1$ & $1$ & $x>1$ \\
$x+1$ & $-$ & $0$ & $+$ & $+$ & $+$ \\
$x-1$ & $-$ & $-$ & $-$ & $0$ & $+$ \\
$f'(x)$ & $+$ & $0$ & $-$ & $0$ & $+$ \\
\end{tblr}
\end{bkcenter}

\noindent たとえば $x<-1$ では、$(x+1)$ も $(x-1)$ も負なので、$3\times(-)\times(-)=(+)$ です。数を入れて確かめると、$f'(-2)=3\times(-1)\times(-3)=9>0$、$f'(0)=3\times1\times(-1)=-3<0$、$f'(2)=3\times3\times1=9>0$ で、表と合っています。

\noindent\textbf{手順3　表にまとめ、$f(x)$ の値を入れる}\quad 上がる所は $\nearrow$、下がる所は $\searrow$ と書きます。$f(-1)=(-1)^3-3\times(-1)=-1+3=2$、$f(1)=1-3=-2$ です。このような表を\textbf{増減表}といいます。

\begin{bkcenter}
\begin{tblr}{colspec={Q[c,m,wd=14mm] *{5}{Q[c,m,wd=16mm]}}, hlines={0.6pt,bkLine}, vline{2}={0.6pt,bkLine}, rowsep=2pt, row{1}={bg=bkTheme!10}}
$x$ & $\cdots$ & $-1$ & $\cdots$ & $1$ & $\cdots$ \\
$f'(x)$ & $+$ & $0$ & $-$ & $0$ & $+$ \\
$f(x)$ & $\nearrow$ & {{\bfseries 2}\\{\scriptsize 極大}} & $\searrow$ & {{\bfseries $-2$}\\{\scriptsize 極小}} & $\nearrow$ \\
\end{tblr}
\end{bkcenter}

\noindent\begin{minipage}[c]{0.42\linewidth}\centering
\begin{tikzpicture}
\begin{axis}[bkplot, width=54mm, height=44mm, xmin=-2.3, xmax=2.5, ymin=-3, ymax=3.4, xtick={-1,1}, ytick={-2,2}, xlabel={$x$}, ylabel={$y$}]
  \addplot[bkcurve, domain=-2.1:2.1] {x^3-3*x};
  \addplot[bkcurve2, line width=1pt, domain=-1.5:-0.6] {2};
  \addplot[bkcurve2, line width=1pt, domain=0.6:1.5] {-2};
  \addplot[only marks, mark=*, mark size=1.8pt, bkInk] coordinates {(-1,2) (1,-2)};
\end{axis}
\end{tikzpicture}
\end{minipage}\hfill
\begin{minipage}[c]{0.55\linewidth}\raggedright
$x=-1$ で、上がるのから下がるのに変わるので、そこが山のてっぺんです。これを\textbf{極大}（きょくだい）といい、そのときの値 $f(-1)=2$ を\textbf{極大値}といいます。$x=1$ は谷の底で\textbf{極小}（きょくしょう）、$f(1)=-2$ が\textbf{極小値}です。極大・極小では、接線が横向き（傾き0、オレンジ）になっています。\bkidx{きょくだい・きょくしょう}{極大・極小}
\end{minipage}

\medskip
\noindent\textbf{「極大」は「いちばん大きい」ではない}\quad 極大値2は、$x=-1$ の\textbf{近くの中で}いちばん大きい値です。グラフ全体では、右へ行くと $f(3)=27-9=18$ のように、2よりずっと大きくなります。山のてっぺんが、全体でいちばん高いとはかぎらないのです。

\bksubsection{3}{範囲が決まっているときの最大と最小}

\label{bk:5-7:saidai}箱の問題では、$x$ は $0<x<6$ の数しか使えませんでした。このように範囲が決まっているときは、極大・極小のほかに、\textbf{範囲のはしの値}もくらべます。山のてっぺんより、はしのほうが高いことがあるからです。

\noindent\begin{minipage}[c]{0.52\linewidth}\raggedright
\textbf{例}\quad $f(x)=x^3-3x$ を、$0\leqq x\leqq3$ の範囲で考えます（$\leqq$ は「小さいか等しい」。2-2）。この範囲にある $f'(x)=0$ の点は $x=1$ だけです（$x=-1$ は範囲の外）。はしの $x=0$、$x=3$ と、$x=1$ の値をくらべます。
\[ f(0)=0,\quad f(1)=-2,\quad f(3)=18 \]
最大値は、はしの $x=3$ のときの18、最小値は $x=1$ のときの $-2$ です。\bkidx{さいだいち}{最大値}\bkidx{さいしょうち}{最小値}
\end{minipage}\hfill
\begin{minipage}[c]{0.44\linewidth}\centering
\begin{tikzpicture}
\begin{axis}[bkplot, width=54mm, height=56mm, xmin=-0.4, xmax=3.6, ymin=-3.5, ymax=21, xtick={1,3}, ytick={-2,18}, xlabel={$x$}, ylabel={$y$}]
  \addplot[bkGray, line width=0.8pt, dashed, domain=-0.3:0] {x^3-3*x};
  \addplot[bkcurve, domain=0:3] {x^3-3*x};
  \addplot[only marks, mark=*, mark size=1.8pt, bkRed] coordinates {(3,18)};
  \addplot[only marks, mark=*, mark size=1.8pt, bkBlue] coordinates {(1,-2) (0,0)};
  \draw[bkGray, dashed, line width=0.5pt] (axis cs:0,18) -- (axis cs:3,18);
  \node[font=\sffamily\scriptsize, anchor=east, text=bkRed!80!black] at (axis cs:2.85,18) {最大（はし）};
  \node[font=\sffamily\scriptsize, anchor=north, text=bkBlue] at (axis cs:1,-2.4) {最小};
\end{axis}
\end{tikzpicture}
\end{minipage}

\bksubsection{4}{はてな？の箱を解く}

\label{bk:5-7:hako}\bkhatenaref の箱の体積は $V=x(12-2x)^2$（$0<x<6$）でした。展開してから微分します。

\noindent\textbf{手順1　展開する}\quad $(12-2x)^2=144-48x+4x^2$（3-8）。これに $x$ をかけます。
\[ V=x(144-48x+4x^2)=144x-48x^2+4x^3 \]
\textbf{手順2　微分する}\quad $(x^n)'=nx^{n-1}$ を1つずつ使います（5-4）。
\[ V'=144-96x+12x^2=12x^2-96x+144 \]
\textbf{手順3　$V'=0$ を解く}\quad どの項も12でわれるので、12でくくってから因数分解します（3-10）。
\[ V'=12(x^2-8x+12)=12(x-2)(x-6) \]
$V'=0$ になるのは $x=2$ と $x=6$。でも、$x=6$ は範囲 $0<x<6$ に入りません（底の1辺が0で、箱にならない）。だから、調べるのは $x=2$ だけです。

\noindent\textbf{手順4　増減表}\quad $x=1$ では $V'=12\times(-1)\times(-5)=60>0$、$x=3$ では $V'=12\times1\times(-3)=-36<0$。2の前で $+$、後で $-$ です。

\begin{bkcenter}
\begin{tblr}{colspec={Q[c,m,wd=14mm] *{5}{Q[c,m,wd=14mm]}}, hlines={0.6pt,bkLine}, vline{2}={0.6pt,bkLine}, rowsep=2pt, row{1}={bg=bkTheme!10}}
$x$ & $0$ & $\cdots$ & $2$ & $\cdots$ & $6$ \\
$V'$ & & $+$ & $0$ & $-$ & \\
$V$ & & $\nearrow$ & {{\bfseries 128}\\{\scriptsize 最大}} & $\searrow$ & \\
\end{tblr}
\end{bkcenter}
\noindent\begin{minipage}[c]{0.4\linewidth}\raggedright
$x=2$ までは体積がふえ、2をこえるとへります。だから、$x=\qty{2}{cm}$ のとき体積は最大で、$2\times8^2=\qty{128}{cm^3}$ です。やってみようの表で見た2が、ちょうどいちばんでした。これが、\bkhatenaref の答えです。
\end{minipage}\hfill
\begin{minipage}[c]{0.57\linewidth}\centering
\begin{tikzpicture}
\begin{axis}[bkplot, width=62mm, height=42mm, xmin=0, xmax=6.4, ymin=0, ymax=150, xtick={1,...,6}, ytick={50,100,128}, xlabel={$x$（cm）}, ylabel={$V$（\unit{cm^3}）}]
  \addplot[bkcurve, domain=0:6] {x*(12-2*x)^2};
  \addplot[only marks, mark=*, mark size=1.6pt, bkInk] coordinates {(1,100) (2,128) (3,108) (4,64) (5,20)};
  \addplot[bkcurve2, line width=1pt, domain=1.2:2.8] {128};
\end{axis}
\end{tikzpicture}
\bkfigcaption{箱の体積。$x=2$ で接線が横向きになる。}
\end{minipage}

\bksubsection{5}{2回微分で、山と谷を見分ける}

\label{bk:5-7:nikai}\noindent\begin{minipage}[c]{0.68\linewidth}\raggedright
山のてっぺんの近くでは、接線の傾きは「正 → 0 → 負」と、だんだん小さくなっていきます。傾き $f'(x)$ がへっていくので、$f'(x)$ を微分した $f''(x)$（2回微分、5-5）は負です。谷の底では反対に、傾きは「負 → 0 → 正」とふえていくので、$f''(x)$ は正です。
\end{minipage}\hfill
\begin{minipage}[c]{0.29\linewidth}\centering
\begin{tikzpicture}[x=0.66mm, y=1mm]
  \draw[bkBlue, line width=1.5pt] (0,0) parabola bend (12,10) (24,0);
  \foreach \xx/\sl in {4/1.1,12/0,20/-1.1} \draw[bkOrange, line width=0.8pt] ({\xx-3},{10-(\xx-12)^2*10/144-3*\sl}) -- ({\xx+3},{10-(\xx-12)^2*10/144+3*\sl});
  \node[font=\sffamily\scriptsize, align=center] at (12,-6) {傾き\\正→0→負\\$f''<0$};
  \draw[bkBlue, line width=1.5pt] (30,10) parabola bend (42,0) (54,10);
  \node[font=\sffamily\scriptsize, align=center] at (42,-6) {傾き\\負→0→正\\$f''>0$};
\end{tikzpicture}
\end{minipage}

\medskip
\noindent\textbf{数で確かめる}\quad $f(x)=x^3-3x$ の山 $x=-1$ の近くで、傾き $f'(x)=3x^2-3$ を計算すると、$f'(-1.5)=3.75$、$f'(-1)=0$、$f'(-0.5)=-2.25$ と、へっていきます。

\begin{itemize}[leftmargin=1.6\zw, itemsep=1pt, topsep=2pt, label={\color{bkTheme}●}]
\item $f'(a)=0$、$f''(a)<0$ なら、$x=a$ は極大（山）
\item $f'(a)=0$、$f''(a)>0$ なら、$x=a$ は極小（谷）
\end{itemize}
$f(x)=x^3-3x$ なら $f''(x)=6x$。$f''(-1)=-6<0$ なので $x=-1$ は極大、$f''(1)=6>0$ なので $x=1$ は極小。増減表と同じ結果です。投げ上げたボールでは、加速度 $\ddot{x}=-10<0$ なので、$\dot{x}=0$ の点はいつも山（最高点）になります。

\bksubsection{6}{「いちばん」を求める手順}

\label{bk:5-7:tejun}文章の問題で最大・最小を求めるときは、次の順に進めます。

\begin{bkfigure}
\begin{tikzpicture}[x=1mm, y=1mm]
  \node[bk node, text width=34mm] (t1) at (0,12) {\textbf{① 式にする}\\大きく（小さく）したい\\量を、$x$ の式で表す};
  \node[bk node, text width=34mm] (t2) at (46,12) {\textbf{② 範囲を決める}\\$x$ が使える範囲\\（長さは正、など）};
  \node[bk node, text width=34mm] (t3) at (92,12) {\textbf{③ 微分して $=0$}\\$f'(x)=0$ を解き、\\範囲の中の解をえらぶ};
  \node[bk node, text width=34mm] (t4) at (23,-14) {\textbf{④ 増減表}\\$f'$ の符号を調べ、\\上がる・下がるを書く};
  \node[bk node, text width=34mm] (t5) at (69,-14) {\textbf{⑤ くらべる}\\極大・極小と、はしの\\値をくらべて決める};
  \draw[bk link] (t1) -- (t2);
  \draw[bk link] (t2) -- (t3);
  \draw[bk link] (t3.south) -- (t4.north east);
  \draw[bk link] (t4) -- (t5);
\end{tikzpicture}
\bkfigcaption{最大・最小の問題の手順}
\end{bkfigure}

\begin{taisetsu}
$f'(x)>0$ なら増加、$f'(x)<0$ なら減少。\\
$f'(x)=0$ の前後で符号が変われば、そこが極大・極小。\\
範囲が決まっているときは、はしの値もくらべて最大・最小を決める。
\end{taisetsu}

\begin{kotoba}
\small
\bkword{増減表}{$f'(x)$ の符号と、$f(x)$ が上がるか下がるかをまとめた表。}
\bkword{極大・極小}{増加から減少に変わる所（山）と、減少から増加に変わる所（谷）。その近くでいちばん大きい・小さい所。}
\bkword{最大値・最小値}{決められた範囲の中で、いちばん大きい値・いちばん小さい値。}
\end{kotoba}

\begin{column}{表計算で「いちばん」をさがす}
式が複雑で $f'(x)=0$ が解けないときは、表計算で $x$ を細かくふって値をくらべます。箱の体積なら、$x$ を0.1きざみにしても0.01きざみにしても、$x=2$ で最大だとわかります。工学では、このように計算機で最大・最小をさがす方法（最適化）がよく使われます。そのもとにあるのも、「傾きが0の所をさがす」考え方です。
\end{column}

%% ---------------- 例題 ----------------
\bksection{reidai}{例題でたしかめよう}

\begin{reidai}[いちばん高い所]
真上に投げたボールの高さが $x=20t-5t^2$（m）のとき、最も高くなる時刻と高さを求めましょう。\par
\kangaekata いちばん高い所では、上がるのから下がるのに変わるので、速度 $\dot{x}$ が0になります（①）。\par
\shiki $\dot{x}=20-5\times2t=20-10t$\par
\hspace*{2.6\zw}$20-10t=0$ より $10t=20$、$t=2$\par
\hspace*{2.6\zw}$t<2$ で $\dot{x}>0$（上がる）、$t>2$ で $\dot{x}<0$（下がる）\par
\hspace*{2.6\zw}$x=20\times2-5\times2^2=40-20=20$\par
\kotae 2秒後に\qty{20}{m}
\end{reidai}

\begin{jibun}
高さが $x=30t-5t^2$（m）のボールの、最も高くなる時刻と高さを求めましょう。
\bkjans{$\dot{x}=30-10t=0$ より3秒後、$90-45=\qty{45}{m}$}
\end{jibun}

\begin{reidai}[増減表と極値]
$f(x)=x^3-6x^2+9x$ の増減表をかき、極大値と極小値を求めましょう。\par
\kangaekata ②の手順1〜3で進めます。\par
\shiki $f'(x)=3x^2-12x+9=3(x^2-4x+3)=3(x-1)(x-3)$\par
\hspace*{2.6\zw}$f'=0$ は $x=1,\,3$\par
\hspace*{2.6\zw}符号：$f'(0)=9>0$、$f'(2)=-3<0$、$f'(4)=9>0$\par
\hspace*{2.6\zw}$f(1)=1-6+9=4$、$f(3)=27-54+27=0$\par
\begin{bkcenter}
\begin{tblr}{colspec={Q[c,m,wd=14mm] *{5}{Q[c,m,wd=13mm]}}, hlines={0.6pt,bkLine}, vline{2}={0.6pt,bkLine}, rowsep=1.5pt, row{1}={bg=bkTheme!10}}
$x$ & $\cdots$ & $1$ & $\cdots$ & $3$ & $\cdots$ \\
$f'(x)$ & $+$ & $0$ & $-$ & $0$ & $+$ \\
$f(x)$ & $\nearrow$ & {{\bfseries 4}\\{\scriptsize 極大}} & $\searrow$ & {{\bfseries 0}\\{\scriptsize 極小}} & $\nearrow$ \\
\end{tblr}
\end{bkcenter}
\kotae $x=1$ で極大値4、$x=3$ で極小値0
\end{reidai}

\begin{jibun}
$f(x)=x^3-12x$ の増減表をかき、極大値と極小値を求めましょう。
\bkjans{$f'=3x^2-12=3(x+2)(x-2)$。$x=-2$ で極大値16、$x=2$ で極小値 $-16$}
\end{jibun}

\begin{reidai}[範囲のある最大・最小（2次の式）]
$f(x)=x^2-4x+5$ の、$0\leqq x\leqq5$ での最大値と最小値を求めましょう。\par
\kangaekata 範囲の中の $f'(x)=0$ の点と、はしの $x=0$、$x=5$ の値をくらべます（③）。\par
\shiki $f'(x)=2x-4=0$ より $x=2$（範囲の中）\par
\hspace*{2.6\zw}$x<2$ で $f'<0$、$x>2$ で $f'>0$ なので、$x=2$ は谷\par
\hspace*{2.6\zw}$f(0)=5$、$f(2)=4-8+5=1$、$f(5)=25-20+5=10$\par
\begin{bkfigure}
\begin{tikzpicture}
\begin{axis}[bkplot, width=62mm, height=44mm, xmin=-0.3, xmax=5.8, ymin=0, ymax=12, xtick={2,5}, ytick={1,5,10}, xlabel={$x$}, ylabel={$y$}]
  \addplot[bkcurve, domain=0:5] {x^2-4*x+5};
  \addplot[only marks, mark=*, mark size=1.8pt, bkRed] coordinates {(5,10)};
  \addplot[only marks, mark=*, mark size=1.8pt, bkBlue] coordinates {(2,1)};
  \addplot[only marks, mark=*, mark size=1.5pt, bkInk] coordinates {(0,5)};
  \node[font=\sffamily\scriptsize, anchor=east, text=bkRed!80!black] at (axis cs:4.8,10.6) {最大（はし）};
  \node[font=\sffamily\scriptsize, anchor=north, text=bkBlue] at (axis cs:2,0.8) {最小（谷）};
\end{axis}
\end{tikzpicture}
\end{bkfigure}
\kotae 最大値10（$x=5$）、最小値1（$x=2$）
\end{reidai}

\begin{jibun}
$f(x)=x^2-6x+10$ の、$0\leqq x\leqq5$ での最大値と最小値を求めましょう。
\bkjans{$f'=2x-6=0$ より $x=3$。$f(0)=10$、$f(3)=1$、$f(5)=5$。最大値10（$x=0$）、最小値1（$x=3$）}
\end{jibun}

\begin{reidai}[範囲のある最大・最小（3次の式）]
$f(x)=x^3-3x^2-9x$ の、$-2\leqq x\leqq4$ での最大値と最小値を求めましょう。\par
\kangaekata 極大・極小が2つとも範囲に入るときは、それと、はしの2つの、合わせて4つの値をくらべます。\par
\shiki $f'(x)=3x^2-6x-9=3(x^2-2x-3)=3(x+1)(x-3)$\par
\hspace*{2.6\zw}$f'=0$ は $x=-1,\,3$（どちらも範囲の中）\par
\hspace*{2.6\zw}$f(-2)=-8-12+18=-2$\par
\hspace*{2.6\zw}$f(-1)=-1-3+9=5$（極大）\par
\hspace*{2.6\zw}$f(3)=27-27-27=-27$（極小）\par
\hspace*{2.6\zw}$f(4)=64-48-36=-20$\par
\begin{bkfigure}
\begin{tikzpicture}
\begin{axis}[bkplot, width=70mm, height=48mm, xmin=-2.5, xmax=4.8, ymin=-31, ymax=9, xtick={-1,3,4}, ytick={-27,5}, xlabel={$x$}, ylabel={$y$}]
  \addplot[bkcurve, domain=-2:4] {x^3-3*x^2-9*x};
  \addplot[only marks, mark=*, mark size=1.8pt, bkRed] coordinates {(-1,5)};
  \addplot[only marks, mark=*, mark size=1.8pt, bkBlue] coordinates {(3,-27)};
  \addplot[only marks, mark=*, mark size=1.5pt, bkInk] coordinates {(-2,-2) (4,-20)};
  \node[font=\sffamily\scriptsize, anchor=south, text=bkRed!80!black] at (axis cs:-1,5.6) {最大};
  \node[font=\sffamily\scriptsize, anchor=west, text=bkBlue] at (axis cs:3.2,-28) {最小};
\end{axis}
\end{tikzpicture}
\end{bkfigure}
\kotae 最大値5（$x=-1$）、最小値 $-27$（$x=3$）
\end{reidai}

\begin{jibun}
$f(x)=x^3-12x$ の、$-3\leqq x\leqq3$ での最大値と最小値を求めましょう。
\bkjans{$f(-3)=9$、$f(-2)=16$、$f(2)=-16$、$f(3)=-9$。最大値16（$x=-2$）、最小値 $-16$（$x=2$）}
\end{jibun}

\begin{reidai}[囲む面積を最大にする]
（1）長さ\qty{40}{m}のロープで長方形を囲みます。面積がいちばん大きくなるのは、たてと横が何mのときでしょう。\par
（2）まっすぐなかべを1辺にして、長さ\qty{40}{m}のロープで残りの3辺を囲みます。かべと直角な辺を $x$ m として、面積が最大になる $x$ と面積を求めましょう。\par
\begin{bkfigure}
\begin{tikzpicture}[x=1mm, y=1mm]
  \fill[bkGreen!15] (0,0) rectangle (30,20);
  \draw[bkOrange!80!black, line width=1.2pt] (0,0) rectangle (30,20);
  \node[font=\sffamily\footnotesize, anchor=east] at (-1,10) {$x$};
  \node[font=\sffamily\footnotesize, anchor=north] at (15,-1) {$20-x$};
  \node[font=\sffamily\footnotesize, anchor=south] at (15,21) {（1）};
  \begin{scope}[shift={(55,0)}]
    \fill[bkGreen!15] (0,0) rectangle (40,20);
    \fill[pattern=north east lines, pattern color=bkGray] (-1,20) rectangle (41,23);
    \draw[bkInk, line width=1.2pt] (-1,20) -- (41,20);
    \draw[bkOrange!80!black, line width=1.2pt] (0,20) -- (0,0) -- (40,0) -- (40,20);
    \node[font=\sffamily\footnotesize, anchor=east] at (-1,10) {$x$};
    \node[font=\sffamily\footnotesize, anchor=north] at (20,-1) {$40-2x$};
    \node[font=\sffamily\footnotesize, anchor=south] at (20,24) {（2）かべ};
  \end{scope}
\end{tikzpicture}
\end{bkfigure}
\kangaekata（1）たて＋横は、ロープの半分の\qty{20}{m}。横は $20-x$。（2）ロープは $x$ が2本と横が1本なので、横は $40-2x$。\par
\shiki（1）$S=x(20-x)=20x-x^2$、$S'=20-2x=0$ より $x=10$。前後で $+$ から $-$\par
\hspace*{2.6\zw}（2）$S=x(40-2x)=40x-2x^2$、$S'=40-4x=0$ より $x=10$。前後で $+$ から $-$\par
\kotae（1）たて・横とも\qty{10}{m}（正方形）、\qty{100}{m^2}\par
\hspace*{2.6\zw}（2）$x=\qty{10}{m}$（横\qty{20}{m}）、\qty{200}{m^2}
\end{reidai}

\begin{jibun}
（1）長さ\qty{24}{m}のロープで長方形を囲みます。面積が最大になるたて・横と面積を求めましょう。\par
（2）かべを1辺にして、長さ\qty{24}{m}のロープで残りの3辺を囲みます。かべと直角な辺を $x$ として、面積が最大になる $x$ と面積を求めましょう。
\bkjans{（1）\nobreak $S=x(12-x)$、$x=6$。たて・横とも\qty{6}{m}、\qty{36}{m^2}\quad（2）\nobreak $S=x(24-2x)$、$S'=24-4x=0$ より $x=\qty{6}{m}$、$6\times12=\qty{72}{m^2}$}
\end{jibun}

\begin{reidai}[箱の体積の最大]
1辺\qty{18}{cm}の正方形の紙の4すみから1辺 $x$ の正方形を切りとって、ふたのない箱をつくります。体積が最大になる $x$ と体積を求めましょう。\par
\kangaekata $V=x(18-2x)^2$（$0<x<9$）。展開してもできますが、積の微分と合成関数の微分（5-6）を使うと、展開しなくてすみます。$x$ と $(18-2x)^2$ のかけ算と見ます。\par
\shiki $(x)'=1$、$\{(18-2x)^2\}'=2(18-2x)\times(-2)$\par
\hspace*{2.6\zw}$V'=1\times(18-2x)^2+x\times2(18-2x)\times(-2)$\par
\hspace*{2.6\zw}$\phantom{V'}=(18-2x)^2-4x(18-2x)$\par
\hspace*{2.6\zw}$\phantom{V'}=(18-2x)\{(18-2x)-4x\}=(18-2x)(18-6x)$\par
\hspace*{2.6\zw}$0<x<9$ で $V'=0$ は $x=3$（$18-2x=0$ の $x=9$ は範囲外）\par
\hspace*{2.6\zw}$x=2$ で $V'=14\times6>0$、$x=4$ で $V'=10\times(-6)<0$\par
\hspace*{2.6\zw}$V=3\times(18-6)^2=3\times144=432$\par
\kotae $x=\qty{3}{cm}$、\qty{432}{cm^3}
\end{reidai}

\begin{jibun}
1辺\qty{24}{cm}の正方形の紙で、同じように箱をつくります。体積が最大になる $x$ と体積を求めましょう。
\bkjans{$V'=(24-2x)^2-4x(24-2x)=(24-2x)(24-6x)$ より $x=\qty{4}{cm}$、$4\times16^2=\qty{1024}{cm^3}$}
\end{jibun}

\begin{reidai}[材料がいちばん少ない缶]
体積が $54\pi$ \unit{cm^3} の、ふたのある円柱の缶をつくります。表面積が最小になる半径 $r$ と高さを求めましょう。\par
\noindent\begin{minipage}[c]{0.5\linewidth}\raggedright
\kangaekata 缶を切り開くと、右の図のように、円が2まい（ふたと底）と、長方形が1まい（側面）になります。長方形の横は、円のまわりの長さ $2\pi r$ です。
\end{minipage}\hfill
\begin{minipage}[c]{0.46\linewidth}\centering
\begin{tikzpicture}[x=1mm, y=1mm]
  \fill[bkBlue!20, draw=bkInk, line width=0.5pt] (0,0) rectangle (44,16);
  \fill[bkYellow!40, draw=bkInk, line width=0.5pt] (22,23) circle (7);
  \fill[bkYellow!40, draw=bkInk, line width=0.5pt] (22,-7) circle (7);
  \draw[bkInk, line width=0.5pt] (22,23) -- (29,23);
  \node[font=\sffamily\scriptsize, anchor=south] at (25.5,23) {$r$};
  \node[font=\sffamily\footnotesize, align=center] at (22,8) {側面（横 $2\pi r$）};
  \node[font=\sffamily\footnotesize, anchor=west] at (44.5,8) {$h$};
\end{tikzpicture}
\end{minipage}\par
\hspace*{2.6\zw}体積 $\pi r^2\times h=54\pi$ より、高さは $h=\dfrac{54}{r^2}$。表面積は\par
\hspace*{2.6\zw}$S=2\times\pi r^2+2\pi r\times h=2\pi r^2+2\pi r\times\dfrac{54}{r^2}=2\pi r^2+\dfrac{108\pi}{r}$\par
\shiki $\dfrac{108\pi}{r}=108\pi r^{-1}$ の微分は $-108\pi r^{-2}$（5-4）なので\par
\hspace*{2.6\zw}$S'=4\pi r-\dfrac{108\pi}{r^2}=0$\par
\hspace*{2.6\zw}両辺に $r^2$ をかけて $4\pi r^3=108\pi$、$r^3=27$、$r=3$\par
\hspace*{2.6\zw}$r=2$ で $S'<0$、$r=4$ で $S'>0$ なので最小。$h=\dfrac{54}{9}=6$\par
\kotae 半径\qty{3}{cm}、高さ\qty{6}{cm}（高さが直径と等しい）
\end{reidai}

\begin{jibun}
体積が $128\pi$ \unit{cm^3} の、ふたのある円柱の缶で、表面積が最小になる半径と高さを求めましょう。
\bkjans{$h=\dfrac{128}{r^2}$、$S=2\pi r^2+\dfrac{256\pi}{r}$。$S'=4\pi r-\dfrac{256\pi}{r^2}=0$ より $r^3=64$、$r=\qty{4}{cm}$、高さ\qty{8}{cm}}
\end{jibun}

\begin{reidai}[電池から取り出す電力を最大にする]
電圧\qty{12}{V}の電池の中には、\qty{2}{\ohm}の抵抗（内部抵抗）があります。この電池に抵抗 $R$ の器具をつなぎます。器具で使われる電力 $P$ が最大になる $R$ と、そのときの電力を求めましょう。\par
\noindent\begin{minipage}[c]{0.56\linewidth}\raggedright
\kangaekata 電流 $I$ は、内部抵抗と $R$ の両方を通ります。直列につないだ抵抗はたし算でまとめられるので、全体の抵抗は $R+2$。$I=\dfrac{V}{R}$（2-15）より、$I=\dfrac{12}{R+2}$。電力は「電圧×電流」で、器具にかかる電圧は $RI$（2-15）なので、
\end{minipage}\hfill
\begin{minipage}[c]{0.4\linewidth}\centering
\begin{tikzpicture}[x=1mm, y=1mm]
  \draw[bkGray, dashed, line width=0.5pt, rounded corners=1.5mm] (-6,-6) rectangle (6,20);
  \node[font=\sffamily\scriptsize, text=bkGray, anchor=north] at (0,-12) {電池};
  \draw[bkInk, line width=0.8pt] (0,-6) -- (0,1) (0,4) -- (0,8) (0,15) -- (0,24) -- (36,24) -- (36,15) (36,3) -- (36,-10) -- (0,-10) -- (0,-6);
  \draw[bkInk, line width=1.2pt] (-3.5,1) -- (3.5,1);
  \draw[bkInk, line width=1.2pt] (-1.8,4) -- (1.8,4);
  \fill[white, draw=bkInk, line width=0.8pt] (-2,8) rectangle (2,15);
  \fill[bkOrange!35, draw=bkInk, line width=0.8pt] (33.5,3) rectangle (38.5,15);
  \node[font=\sffamily\scriptsize, anchor=west] at (7,11.5) {\qty{2}{\ohm}};
  \node[font=\sffamily\scriptsize, anchor=west] at (7,2.5) {\qty{12}{V}};
  \node[font=\sffamily\footnotesize, anchor=west] at (39.5,9) {$R$};
  \draw[bk arrow, bkRed, line width=0.9pt] (14,26.5) -- (24,26.5);
  \node[font=\sffamily\scriptsize, text=bkRed, anchor=south] at (19,27.5) {$I$};
\end{tikzpicture}
\end{minipage}\par
\hspace*{2.6\zw}$P=RI\times I=R\times\dfrac{144}{(R+2)^2}=144R(R+2)^{-2}$\par
\shiki 積の微分と合成関数の微分（5-6）で\par
\hspace*{2.6\zw}$P'=144\{1\times(R+2)^{-2}+R\times(-2)(R+2)^{-3}\}$\par
\hspace*{2.6\zw}$(R+2)^{-2}=(R+2)\times(R+2)^{-3}$ なので、$(R+2)^{-3}$ でくくると\par
\hspace*{2.6\zw}$P'=144(R+2)^{-3}\{(R+2)-2R\}=\dfrac{144(2-R)}{(R+2)^3}$\par
\hspace*{2.6\zw}$R<2$ で $P'>0$、$R>2$ で $P'<0$ なので、$R=2$ で最大\par
\kotae $R=\qty{2}{\ohm}$（内部抵抗と同じ）、$P=\dfrac{144\times2}{4^2}=\dfrac{288}{16}=\qty{18}{W}$
\end{reidai}

\noindent\textbf{表で確かめる}\quad $R$ を変えて $P=\dfrac{144R}{(R+2)^2}$ を計算すると、下の表のように、$R=2$ のときがいちばん大きくなっています。抵抗が小さすぎると、電流は大きくても器具の電圧が小さく、抵抗が大きすぎると、電流が小さくなってしまうのです。
\begin{bkcenter}
\begin{bktable}{colspec={Q[c,m,wd=16mm] *{5}{Q[c,m,wd=12mm]}}, rowsep=1.5pt}
$R$（\unit{\ohm}） & 1 & 2 & 3 & 4 & 8 \\
$P$（W） & 16 & 18 & 17.28 & 16 & 11.52 \\
\end{bktable}
\end{bkcenter}

\begin{jibun}
電圧\qty{6}{V}、内部抵抗\qty{3}{\ohm}の電池に抵抗 $R$ をつなぎます。$R$ で使われる電力 $P$ を $R$ の式で表し、$P$ が最大になる $R$ と電力を求めましょう。
\bkjans{$I=\dfrac{6}{R+3}$、$P=RI^2=\dfrac{36R}{(R+3)^2}$。$P'=\dfrac{36(3-R)}{(R+3)^3}$ より $R=\qty{3}{\ohm}$、$P=\dfrac{108}{36}=\qty{3}{W}$}
\end{jibun}

%% ---------------- まちがえやすいところ・物理工学 ----------------
\bksection{tsumazuki}{まちがえやすいところをチェック}

\begin{tsumazuki}{$f'(x)=0$ なら、いつも極大か極小？}
\bkng{$f(x)=x^3$ で $f'(0)=0$ だから、$x=0$ は極大か極小}
\bkok{$x=0$ の前後で $f'>0$ のまま。極大でも極小でもない}
\vspace{2pt}\noindent\begin{minipage}[c]{0.76\linewidth}\raggedright
\strut$f(x)=x^3$ では $f'(x)=3x^2$ なので、$f'(0)=0$ です。でも、$x=0$ の前でも後でも $f'>0$ で、ずっと上がり続けています（とちゅうで一瞬だけ平らになる）。$f'(x)=0$ でも、前後で符号が変わらなければ、極大でも極小でもありません。増減表をかいて、符号が変わるかを確かめましょう。
\end{minipage}\hfill
\begin{minipage}[c]{0.21\linewidth}\centering
\begin{tikzpicture}
\begin{axis}[bkplot, width=28mm, height=30mm, xmin=-1.5, xmax=1.6, ymin=-2, ymax=2.2, xtick=\empty, ytick=\empty, xlabel={$x$}, ylabel={$y$}]
  \addplot[bkcurve, domain=-1.25:1.25] {x^3};
  \addplot[bkcurve2, line width=1pt, domain=-0.6:0.6] {0};
\end{axis}
\end{tikzpicture}
\end{minipage}
\end{tsumazuki}

\begin{tsumazuki}{範囲のはしを調べ忘れる}
\bkng{$f(x)=x^3-3x$（$-2\leqq x\leqq3$）の最大値は、極大値の2}
\bkok{はしの $f(3)=27-9=18$ もくらべて、最大値は18}
\noindent\begin{minipage}[c]{0.62\linewidth}\raggedright
範囲が決まっているとき、極大値がいちばん大きいとはかぎりません（右の図。\bkreidairef{3}でも、はしの $x=5$ が最大）。極大・極小の値と、はしの値を全部くらべます。
\end{minipage}\hfill
\begin{minipage}[c]{0.34\linewidth}\centering
\begin{tikzpicture}
\begin{axis}[bkplot, width=40mm, height=36mm, xmin=-2.4, xmax=3.6, ymin=-4, ymax=21, xtick={-1,3}, ytick={2,18}, xlabel={$x$}, ylabel={$y$}]
  \addplot[bkcurve, line width=1.3pt, domain=-2:3] {x^3-3*x};
  \addplot[only marks, mark=*, mark size=1.6pt, bkRed] coordinates {(3,18)};
  \addplot[only marks, mark=*, mark size=1.4pt, bkInk] coordinates {(-1,2)};
\end{axis}
\end{tikzpicture}
\end{minipage}
\end{tsumazuki}

\begin{tsumazuki}{範囲の外の解も答えにしてしまう}
\bkng{はてな？の箱で、$V'=0$ だから $x=2$ と $x=6$ が答え}
\bkok{$x=6$ は $0<x<6$ に入らない（底の1辺が0）。答えは $x=2$ だけ}
$f'(x)=0$ の解が出たら、問題の範囲に入っているかを、いつも確かめます。とくに長さ・時間・体積は、0やマイナスになると意味がなくなります。
\end{tsumazuki}

\begin{butsuri}
\begin{itemize}[leftmargin=1.2\zw, itemsep=2pt, topsep=0pt, label={\color{bkPurple}●}]
\item \textbf{缶の形}\\
同じ量が入る円柱の缶のうち、材料（表面積）がいちばん少ないのは、高さが直径と等しい形です（\bkreidairef{7}）。ほんとうの缶は、ふたの厚みや、持ちやすさ・積みやすさも考えて形を決めています。
\item \textbf{いちばん多く電気をとり出す}\\
器具の抵抗を電池の内部抵抗と同じにすると、器具で使える電力がいちばん大きくなります（\bkreidairef{8}）。電気工学で「整合」（マッチング）とよばれる考え方で、アンテナやスピーカーをつなぐときに使われます。
\item \textbf{ボールをいちばん遠くへ}\\
空気の抵抗がなければ、地面から45°の向きに投げたとき、いちばん遠くまで届きます。三角関数の微分（5-15）を学ぶと、この章の方法で確かめられます。
\end{itemize}
\end{butsuri}

%% ---------------- 章末問題 ----------------
\bkendproblems
\bkproblemgroup{基本}{3}

\bkproblem $f(x)=x^3-3x^2$ の増減表をかき、極大値と極小値を求めましょう。
\bkpans{$f'=3x^2-6x=3x(x-2)$。$x=0$ で極大値0、$x=2$ で極小値 $-4$}

\bkproblem 高さ $x=40t-5t^2$（m）のボールの、最高点の時刻と高さを求めましょう。
\bkpans{$\dot{x}=40-10t=0$ より4秒後、$160-80=\qty{80}{m}$}

\bkproblem $f(x)=-x^2+6x$ の、$0\leqq x\leqq4$ での最大値と最小値を求めましょう。
\bkpans{$f'=-2x+6$ より $x=3$。$f(0)=0$、$f(3)=9$、$f(4)=8$。最大値9（$x=3$）、最小値0（$x=0$）}

\bkproblem 長さ\qty{60}{cm}の針金で長方形をつくります。面積が最大になるたて・横と、その面積を求めましょう。
\bkpans{$S=x(30-x)$、$x=15$。たて・横とも\qty{15}{cm}、\qty{225}{cm^2}}

\bkproblem $f(x)=2x^3-3x^2-12x$ の極大値と極小値を求めましょう。
\bkpans{$f'=6x^2-6x-12=6(x-2)(x+1)$。$x=-1$ で極大値7、$x=2$ で極小値 $-20$}

\bkproblem $x>0$ で、$f(x)=x+\dfrac{9}{x}$ の最小値を求めましょう。
\bkpans{$f'=1-\dfrac{9}{x^2}=0$ より $x^2=9$、$x>0$ なので $x=3$。最小値 $3+3=6$}

\bkproblemgroup{チャレンジ}{4}

\bkproblem 1辺\qty{30}{cm}の正方形の紙の4すみから1辺 $x$（cm）の正方形を切りとって、ふたのない箱をつくります。体積 $V=x(30-2x)^2$ が最大になる $x$ と、そのときの体積を求めましょう（積の微分と合成関数の微分を使ってもよい）。
\bkpans{$V'=(30-2x)^2-4x(30-2x)=(30-2x)(30-6x)$。$x=\qty{5}{cm}$、$5\times20^2=\qty{2000}{cm^3}$}

\bkproblem 体積が $16\pi$ \unit{cm^3} の円柱の缶をつくります。半径を $r$ とすると、表面積は $S=2\pi r^2+\dfrac{32\pi}{r}$ です。$S$ が最小になる $r$ と、そのときの高さを求めましょう。
\bkpans{$S'=4\pi r-\dfrac{32\pi}{r^2}=0$ より $r^3=8$、$r=\qty{2}{cm}$。高さ $\dfrac{16}{2^2}=\qty{4}{cm}$（直径と等しい）}

\bkproblem 電圧\qty{9}{V}、内部抵抗\qty{1.5}{\ohm}の電池に抵抗 $R$ をつなぐと、$R$ の電力は $P=\dfrac{81R}{(R+1.5)^2}$（W）です。$P$ が最大になる $R$ と電力を求めましょう。
\bkpans{$P'=\dfrac{81(1.5-R)}{(R+1.5)^3}$ より $R=\qty{1.5}{\ohm}$、$P=\dfrac{121.5}{9}=\qty{13.5}{W}$}

\begin{furikaeri}
\begin{checklist}
\item 導関数の符号から、上がる・下がるを判断できる。\hfill{\footnotesize\sffamily\color{bkTheme}→\pageref{bk:5-7:fugou}ページ}
\item 増減表をかいて、極大・極小を求められる。\hfill{\footnotesize\sffamily\color{bkTheme}→\pageref{bk:5-7:zougen}ページ}
\item 範囲のはしもくらべて、最大・最小を求められる。\hfill{\footnotesize\sffamily\color{bkTheme}→\pageref{bk:5-7:saidai}ページ}
\end{checklist}
\end{furikaeri}

\bkfukidashi{mina}{接線は、グラフの「近くのようす」をよく表しているんだね。次の章では、接線を使って、少し先の値を予想するよ。}
